\section{Introduction}\label{sec:introduction}

Simulation optimization is often most valuable precisely where function
evaluation is slow, noisy, and safety constrained.  A policy may contain
hundreds or thousands of coordinated controls, while only a handful of target
simulations are affordable.  Examples include reserve schedules, multi-node
inventory rules, and resource-allocation profiles.  In these settings, the
central statistical difficulty precedes the choice of an acquisition
function: an unstructured initial design has negligible probability of
intersecting a small feasible basin.  A sophisticated sequential optimizer
cannot improve a basin it never observes.

Historical tasks create a possible escape from this difficulty.  Related
systems may have been optimized in earlier markets, networks, or operating
regimes.  Their observations can reveal policy shapes that repeatedly balance
objective quality and chance feasibility.  Yet reusing those observations is
delicate.  Raw coordinates may change dimension, the meaning of individual
controls may not align across tasks, and a target-specific hand-coded anchor
would turn transfer into hidden oracle information.  Moreover, an attractive
recommendation and a deployment-safe recommendation are different statistical
objects.  Reusing the same observations for search and final certification can
make the apparent evaluation budget and the safety claim both misleading.

This paper treats those issues as one information-design problem.  Before a
target task is opened, source outcomes are used to construct a small
\emph{structural proposal atlas}: a set of normalized policy profiles that is
defined in a dimension-equivariant low-frequency coordinate.  Source profiles
are ranked by objective and chance-margin evidence, target-label-free profiles
provide universal structural support, and a maximin rule fills a fixed initial
design.  The atlas is then frozen.  Any target optimizer can consume it; in the
primary implementation we use the standard sparse axis-aligned subspace
Bayesian optimizer (SAASBO).  After search, a separate verifier receives fresh
samples from a frozen, ordered shortlist and never returns those samples to the
optimizer.

The separation is substantive rather than cosmetic.  It lets us ask three
distinct questions.  Does the transferred frontend reach a feasible basin?
Does online optimization improve objective quality once coverage is
established?  Can the final recommendation be independently certified?  It
also forces a transparent cost accounting: source simulation, target search,
and terminal verification are reported separately.

Our contributions are as follows.

\begin{enumerate}[label=(\roman*)]
\item We introduce a target-label-free, dimension-equivariant risk-objective
proposal atlas.  Its low-frequency coordinate has fixed dimension even when
the raw policy dimension changes, and its source-monotone extension is
fail-closed when source directions disagree.

\item We establish the relevant theoretical boundary.  No proper finite atlas
can cover every possible held-out feasible set.  Under an explicit
source-to-target support condition, coordinate error bound, margin regularity,
and safe-depth condition, however, the atlas contains a feasible policy.  The
sufficient condition depends on coordinate geometry rather than nominal policy
dimension.  These and the terminal verification results are machine checked in
Lean 4 without unproved placeholders.

\item We design an optimizer-independent verifier with familywise control of
unsafe deployment and an objective switch guard.  Search, objective-comparison,
and safety-verification samples are tracked separately, so a target search
budget is never reported as a total simulation budget.

\item We perform a causal frontend/backend evaluation.  At dimension 1,000,
the frozen atlas produces true-feasible and independently certified
recommendations in all 60 trials with proposal-only, stacked-GP, and canonical
SAASBO backends.  Matched common-Sobol frontends produce zero, one, and zero,
respectively.  At dimension 10,000, all 30 trials are certified across target
search budgets from 10 to 40 calls.  An electricity-storage holdout confirms a
large total-cost advantage over unstructured Sobol, while a post-confirmatory
structured control shows no advantage over an obvious low-frequency grid.

\item We retain negative findings that delimit the contribution.  Canonical
SAASBO is a replaceable backend rather than a novel algorithm.  A cumulative
heteroscedastic variance model recovers variance shape much better than pooled
variance but does not improve optimization outcomes, so it is reported as a
calibration diagnostic rather than a second headline contribution.
\end{enumerate}

The resulting claim is deliberately narrower than ``solving arbitrary
10,000-dimensional optimization in ten evaluations.''  The method succeeds
when related tasks share a stable, low-dimensional structural coordinate; the
paper states, audits, and stress-tests that condition.  Section
\ref{sec:literature} positions the work.  Sections \ref{sec:problem}--
\ref{sec:theory} give the information contract, atlas, and guarantees.
Sections \ref{sec:experiments} and \ref{sec:results} present the experiments,
and Section \ref{sec:discussion} discusses limits and operational use.
